
\begin{frame}
\frametitle{Solutions - Translation layers - Wine and Proton}
\begin{redbox}
\textbf{Wine} ("Wine Is Not an Emulator")
\vspace{5pt}
\begin{itemize}
    \setlength{\itemsep}{8pt}
    \item Wine is an open source compatibility layer for running Windows software on unix-like systems (MacOS, Linux)
    \item Works based on reverse engineering the Windows kernel
    \item Runs on the fly in user space, using wineserver daemon 
\end{itemize}
\end{redbox}
\begin{yellowbox}
\textbf{Proton}
\vspace{5pt}
\begin{itemize}
    \setlength{\itemsep}{8pt}
    \item Proton is a fork of wine bundled with other software, such as VKD3D, a translation layer for DirectX $\rightarrow$ Vulkan
    \item Proton runs out of the box (using Steam Play)
\end{itemize}
\end{yellowbox}

\end{frame}

\begin{frame}
\frametitle{Wine and Proton - Advantages and Disadvantages}
\begin{theorembox}
\textbf{Advantages}
\begin{itemize}
    \item Minimal input from the developer to port their system using wine or proton
    \item Lower performance overhead as wine/proton acts as a translation layer, and not a full virtual machine
    \item Open source
\end{itemize}
\end{theorembox}
\begin{redbox}
\textbf{Disadvantages}
\begin{itemize}
    \item Apps often don't work out of the box, requiring configuration (improved with proton)
    \item Not all software runs perfectly, if at all
\end{itemize}
\end{redbox}
\end{frame}

\begin{frame}
\frametitle{Wine and Proton - Evaluation}
\begin{theorembox}
\begin{itemize}
    \setlength{\itemsep}{8pt}
    \item \textbf{Stability} - Stability is high, with occasional crashes for misconfigured/incompatible applications
    \item \textbf{Performance} - Marginal effect on performance $<$5\%
    \item \textbf{Features} - Lack of driver support and limited anti-cheat systems as kernel level access cannot be granted.
    \item \textbf{Updates} - Major updates for wine occur roughly every year, Major updates for proton occur every 6-12 months
\end{itemize}
\end{theorembox}
\end{frame}